\def\level{Première}  % Classe à droite
\def\course{NSI}
\def\chaptername{Boucles FOR et WHILE} % Titre au centre
\def\docname{AP04-S}        % Identifiant à gauche

\pagestyle{cours_FP}
\makeheader{} % Affiche le bandeau

%\section{La boucle FOR}


	\begin{easybox}{La boucle FOR en python}

				\begin{minipage}[b]{0.33\linewidth}
	\begin{CodePythontexAlt}[Largeur=1\linewidth,Centre,PremLigne=1,Lignes=false,]{}
for i in range(0, 7):
	print(i, end=" ")
# affichage: 
# 0 1 2 3 4 5 6
	
				\end{CodePythontexAlt}
		\end{minipage}\hfill
		\begin{minipage}[b]{0.33\linewidth}
	\begin{CodePythontexAlt}[Largeur=1\linewidth,Centre,PremLigne=1,Lignes=false,]{}
for i in range(2, 7):
	print(i, end=" ")
# affichage: 
# 2 3 4 5 6
					
				\end{CodePythontexAlt}
		\end{minipage}\hfill
		 \begin{minipage}[b]{0.33\linewidth}
	\begin{CodePythontexAlt}[Largeur=1\linewidth,Centre,PremLigne=1,Lignes=false,]{}
for i in range(5):
	print(i, end=" ")
# affichage: 
# 0 1 2 3 4
	
				\end{CodePythontexAlt}
		\end{minipage}
	\end{easybox}
	
	\begin{easybox}{Itération sur une chaine de caractères}{}
	
	\begin{minipage}{0.48\linewidth}
	\begin{CodePythontexAlt}[Largeur=1\linewidth,Centre,PremLigne=1]{}
	mot = 'Hello'
	for i in range(0, len(mot)):
		print(mot[i])
			
	\end{CodePythontexAlt}
	\end{minipage}\hfill
	\begin{minipage}{0.48\linewidth}
		\begin{CodePythontexAlt}[Largeur=1\linewidth,Centre,PremLigne=1]{}
	mot = 'Hello'
	for lettre in mot:
		print(lettre)	
		\end{CodePythontexAlt}
	
	
	\end{minipage}
	\end{easybox}
		
	
	\begin{easybox}{Algorithmes de Recherche dans un itérable}
	
	\begin{minipage}{0.48\linewidth}
		\begin{CodePythontexAlt}[Largeur=1\linewidth,Centre,PremLigne=1, Lignes=false,]{}
				def est_present(mot, lettre_cherchee):
					for lettre in mot:
						if lettre == lettre_cherchee:
							return True
					return False
				\end{CodePythontexAlt}
		
	\end{minipage}\hfill
	\begin{minipage}{0.48\linewidth}
		\begin{CodePythontexAlt}[Largeur=1\linewidth,Centre,PremLigne=1, Lignes=false,]{}
				def est_present(mot, lettre_cherchee):
					for i in range(len(mot)):
						if mot[i] == lettre_cherchee:
							return True
					return False
				\end{CodePythontexAlt}
		
	\end{minipage}
		
		\end{easybox}

	\begin{easybox}{Algorithmes de comptage dans un itérable}
		
		\begin{minipage}{0.48\linewidth}
			\begin{CodePythontexAlt}[Largeur=1\linewidth,Centre,PremLigne=1, Lignes=false,]{}
			def compte(mot, lettre_cherchee):
				c = 0
				for lettre in mot:
					if lettre == lettre_cherchee:
						c = c + 1
				return c
				
			\end{CodePythontexAlt}
	
		\end{minipage}\hfill
		\begin{minipage}{0.48\linewidth}
		\begin{CodePythontexAlt}[Largeur=1\linewidth,Centre,PremLigne=1, Lignes=false,]{}
			def compte(mot, lettre_cherchee):
				c = 0
				for i in range(len(mot)):
					if mot[i] == lettre_cherchee:
						c = c + 1
				return c
				
			\end{CodePythontexAlt}
	
		\end{minipage}
		

	\end{easybox}

	\begin{easybox}{Algorithmes de somme}
			\begin{CodePythontexAlt}[Largeur=1\linewidth,Centre,PremLigne=1, Lignes=false,]{}
			def somme(n):
				''' Renvoie la somme des entiers de 1 à n inclus'''
				s = 0
				for i in range(1, n+1):
					s = s + i
				return s

			\end{CodePythontexAlt}
			\end{easybox}

			\renewcommand{\arraystretch}{1}


		%\section{La boucle WHILE}


		\begin{easybox}{Prototypage de la boucle WHILE en python}
	\verb|condition1| est un booléen qui est évalué à chaque tour de boucle.  Tant que \verb|condition1| \verb|True|, on effectue les instructions de la boucle. Lorsque \verb|condition1| devient \verb|False|, on sort de la boucle.

	
	\begin{minipage}{0.48\linewidth}
			\begin{CodePythontexAlt}[Largeur=1\linewidth,Centre,PremLigne=1]{}
			n = 0
			u = 1
			while u < 1000:
				n = n + 1
				u = u * 2
			print(n)
			print(u)
		\end{CodePythontexAlt}
				
		\end{minipage}\hfill
		\begin{minipage}{0.48\linewidth}
			\begin{CodePythontexAlt}[Largeur=1\linewidth,Centre,PremLigne=1]{}
			mot = 'abracadabra'
			c = 0
			trouve = False
			while c < len(mot) and not(trouve):
				trouve = mot[c] == 'r'
				c = c + 1
			print(trouve)
			\end{CodePythontexAlt}
		\end{minipage}	

			\end{easybox}

